\documentclass[11pt]{article}

\usepackage[margin=1.05in]{geometry}
\usepackage{amsmath,amssymb,amsthm,array}
\usepackage{url}
\usepackage[hidelinks]{hyperref}

\numberwithin{equation}{section}

\newcommand{\A}{\alpha'}
\newcommand{\bC}{\mathbb C}
\newcommand{\bR}{\mathbb R}
\newcommand{\bZ}{\mathbb Z}
\newcommand{\cB}{\mathcal B}
\newcommand{\cC}{\mathcal C}
\newcommand{\cF}{\mathcal F}
\newcommand{\cG}{\mathcal G}
\newcommand{\cH}{\mathcal H}
\newcommand{\cI}{\mathcal I}
\newcommand{\cL}{\mathcal L}
\newcommand{\cM}{\mathcal M}
\newcommand{\cQ}{\mathcal Q}
\newcommand{\cR}{\mathcal R}
\newcommand{\cS}{\mathcal S}
\newcommand{\cV}{\mathcal V}
\newcommand{\cZ}{\mathcal Z}
\newcommand{\dd}{\mathrm d}
\newcommand{\ee}{\mathrm e}
\newcommand{\ii}{\mathrm i}
\newcommand{\im}{\operatorname{Im}}
\newcommand{\rank}{\operatorname{rank}}
\newcommand{\Aut}{\operatorname{Aut}}
\newcommand{\Span}{\operatorname{span}}
\newcommand{\Ker}{\operatorname{Ker}}
\newcommand{\tr}{\operatorname{tr}}
\newcommand{\vol}{\operatorname{vol}}
\newcommand{\wt}{\widetilde}
\newcommand{\wh}{\widehat}
\newcommand{\vev}[1]{\left\langle #1\right\rangle}
\newcommand{\cA}{\mathcal A}
\newcolumntype{L}[1]{>{\raggedright\arraybackslash}p{#1}}

\newtheorem{definition}{Definition}
\newtheorem{theorem}{Theorem}
\newtheorem{regulatorassumption}{Regulator Assumption}
\newtheorem{remark}{Remark}

\title{Numerical Bosonic String Integrands from One-Face Ribbon Graphs}
\author{GPT 5.5\\under Xi Yin's supervision}
\date{\today}

\begin{document}
\maketitle

\begin{abstract}
This note extracts the numerical construction of bosonic string integrands from
the local string Monte Carlo note, the Markdown implementation notes, and the
Python scripts in the covariant ribbon-graph repository.  The purpose is not to
replace the draft, but to isolate the precise mathematical object computed by
the code, the setup in which it represents a continuum worldsheet CFT quantity,
the one regulator assumption, and the checks that support the construction.

The central object is a Riemann surface obtained by sewing the boundary of a
disk according to a one-face metric ribbon graph.  The matter determinant and
compact-boson winding sum are computed directly from a discretized identity
wavefunctional on the disk.  The holomorphic data, period matrix, prime form,
Riemann constants, and \(bc\) correlators are reconstructed from the same disk
sewing.  These data give a numerical representative of the matter and ghost
factors in the bosonic string measure, up to the standard issues of Weyl-frame
renormalization, moduli-independent normalization, automorphism factors, and the
treatment of ribbon-cell boundaries.  The native moduli space in this note is
\(\cM_{g,1}\); for partition-function applications the marked puncture carries
the identity operator.
\end{abstract}

\tableofcontents

\section{Scope and status}

The sources reviewed for this note are:
\begin{itemize}
\item the local string Monte Carlo note
  \path{covariant formalism/tex/string MC.tex};
\item the Strebel/ribbon working note \path{Strebel.tex} and the draft
  outline in \path{draft/};
\item the Markdown implementation notes
  \path{covariant formalism/markdown/riemann_surface_tools.md} and
  \path{covariant formalism/markdown/genus3_bc_chi18_ratio_check.md};
\item the numerical scripts in \path{covariant formalism/python/}, especially
  the graph-generation, determinant, compact-boson, period-matrix,
  Riemann-surface, scaling-sweep, and Mumford-check scripts listed explicitly
  in Section~11;
\item the relevant sections of the external string notes source
  \url{https://github.com/xiyin137/stringbook/},
  namely conformal gauge with moduli, the holomorphic reformulation of
  \(b\)-ghost insertions, the ghost number anomaly, and the higher-loop
  bosonic string measure.
\end{itemize}

I will use the following status labels.

\begin{description}
\item[Exact CFT.] A statement follows from the continuum CFT or string measure
formalism, independent of the lattice regulator.
\item[Setup.] A statement fixes the class of graphs, coordinates, or
normalization conventions used in this note.
\item[Theorem/geometric input.] A statement is a standard mathematical input,
for example the Strebel-Kontsevich cell decomposition or a continuum CFT
identity.
\item[Regulator assumption.] A statement asserts convergence of the
discretized disk construction to the continuum quantity after allowed local
counterterms.  This is the only nontrivial assumption made by the numerical
construction.
\item[Check.] A statement records a diagnostic already implemented in the
repository.
\end{description}
In the body below these labels appear either as theorem-like environments or as
bold lead-ins to the relevant paragraphs.  They are a bookkeeping device for
separating definitions from continuum facts, numerical evidence, and the single
regulator assumption.

The output of the present construction should be read as a local integrand or
section over the moduli represented by a fixed-perimeter one-face ribbon cell,
namely a cell in \(\cM_{g,1}\).  It is not, by itself, a completed global
integral until the cell measure, automorphism factors, and boundary strata are
treated consistently.

\section{Setup, theorems, and the regulator assumption}

\noindent\textbf{Setup.}
Unless stated otherwise, the numerical formulas use \(\A=1\).  Factors of
\(\A\) are restored in a few continuum formulas to identify the convention.
The worldsheet is a closed oriented genus \(g\) Riemann surface \(\Sigma\).
The point of this section is to separate definitions and standard theorems from
the one real analytic assumption in the numerical work.

\begin{definition}[Ribbon-cell setup]
The primary numerical coordinate system is a positive-length, one-face metric
ribbon graph \(\Gamma\) of genus \(g\).  On the open top-dimensional cells the
graph is trivalent, and hence
\begin{equation}
        E=6g-3,\qquad V=4g-2,\qquad F=1.
\end{equation}
Degenerate cells with vanishing edge lengths are boundary strata and are not
part of the present local construction.
\end{definition}

\begin{theorem}[Strebel-Kontsevich input]
For a genus \(g\) Riemann surface with one marked point and one positive
perimeter, the Strebel differential determines a one-face metric ribbon graph;
conversely the corresponding metric ribbon graphs give the standard cell
decomposition of
\begin{equation}
        \cM_{g,1}\times \bR_{+},
\end{equation}
with the usual orbifold identifications by graph automorphisms.
\end{theorem}

\begin{remark}[Marked point and perimeter]
The preceding theorem is not an assumption of the numerical method; it is the
geometric setup.  In this note we take \(\cM_{g,1}\), with fixed Strebel
perimeter, as the native moduli space.  For partition-function applications the
marked puncture is assigned the identity operator.  Thus the marked point is
part of the data being integrated over, rather than something to be silently
removed.  If one later wants an unmarked vacuum amplitude on \(\cM_g\), one
must specify a separate pushforward or fiber-normalization convention.
\end{remark}

\begin{definition}[Weyl frame, zero modes, and winding sector]
The lattice regulator is defined in the disk frame associated with the Strebel
differential.  Individual matter or ghost partition functions are therefore
Weyl-frame dependent, while Weyl-invariant combinations, ratios at fixed
central charge, and the critical matter-plus-ghost combination are the intended
continuum quantities.

For a noncompact boson, the constant zero mode is factored as a target-space
volume and the determinant is taken on the orthogonal complement of constants.
For a compact boson \(X\sim X+2\pi R\), the constant zero mode contributes
\(2\pi R\).  Winding around the marked puncture is set to zero; equivalently,
only edge-shift configurations satisfying the vertex incidence constraints are
summed.  These are conventions defining the regulated quantity.
\end{definition}

\begin{definition}[Holomorphic data conventions]
Period matrices are computed from \(A\)-normalized holomorphic one-forms.  Abel
maps use a fixed basepoint in the disk, and physical formulas only use
basepoint-independent combinations.  Prime forms use an odd theta
characteristic and principal square-root branches; the independence of the
prime form from the odd characteristic is a standard theorem and is checked
numerically.

For \(g\geq2\), the note and the current checks use the canonical construction
of the Riemann constant vector described in the implementation notes: zeros of
an \(A\)-normalized holomorphic one-form determine a canonical divisor, and
half-period candidates are selected by Riemann-vanishing tests.  The
higher-genus \(\sigma\) normalization is fixed from the \(\lambda=1\)
Verlinde-Verlinde identity together with the chiral determinant convention.
The older direct cycle-integral routine is not used as an input to the claims
in this note.
\end{definition}

\begin{definition}[Normalization convention]
All absolute amplitudes in this note are defined only up to constants that are
independent of moduli and of the compactification radius.  The implemented
checks compare ratios or local sections for which these constants cancel.
\end{definition}

\begin{regulatorassumption}[Continuum limit modulo local counterterms]
For integer edge lengths, set
\begin{equation}
        m=\sum_{e=1}^{E}\ell_e,\qquad L=2m.
\end{equation}
The continuum limit is taken by scaling all edge lengths by a common large
integer while keeping their ratios fixed.  The nontrivial analytic assumption
is that the resulting discretized disk path integral converges to the intended
continuum CFT quantity after multiplication by local counterterms allowed by
the symmetries of the regulator, such as a lattice cosmological term
proportional to \(L\) and an Euler-curvature term.  The remaining freedom is a
finite moduli-independent normalization.

The large-\(L\) fits used in the code are the practical implementation of this
assumption.  They require all edge lengths to be large enough that short-edge
artifacts are negligible; the implementation notes use minimum edge lengths of
order \(200\) for stable fits.
\end{regulatorassumption}

\begin{remark}
There is a counting-convention issue in the working files.  The draft outline
cites the Bacher--Vdovina counts \(9\) and \(1726\) for one-vertex genus-two
and genus-three triangulations.  The current compact-boson file
\path{compact_partition.py} stores \(9\) genus-two one-face graphs, and the
current genus-three data file \path{genus3_ribbon_graph_data.txt} declares
\(1726\) stored one-face genus-three ribbon graphs.  The current generator
describes its output as isomorphism-reduced under graph automorphisms.  An
older sentence in \path{string MC.tex} says that the compact-boson T-duality
check was run over \(4\) genus-two topologies and \(818\) genus-three
topologies.  The reviewed repo files do not identify the enumeration convention
behind the older \(4\) and \(818\) numbers.  Until this is audited, the
current-code counts should be quoted as \(9\) and \(1726\), and the older
\(4\) and \(818\) statements should be treated as unresolved.  No formula below
depends on these counts.
\end{remark}

\section{Continuum string measure}

\noindent\textbf{Exact CFT.}
This section records the continuum target of the numerical construction.  In
conformal gauge, after fixing puncture positions and a fiducial family of
metrics \(\hat g(t)\), the Faddeev-Popov procedure gives the usual insertion
\begin{equation}
        \cB_{t^k}
        =
        \frac{1}{4\pi}
        \int_{\Sigma}\dd^2\sigma\,\sqrt{\hat g}\,
        b^{ab}\frac{\partial \hat g_{ab}}{\partial t^k}.
        \label{eq:Bmetric}
\end{equation}
Thus a genus \(h\), \(n\)-point bosonic string amplitude has the schematic
form
\begin{equation}
        \cA_h[V_1,\ldots,V_n]
        =
        N_{h,n}\int_{\cM_{h,n}}\Omega_{6h-6+2n},
        \qquad
        \Omega=\left\langle \ee^{\cB}\prod_i\cV_i\right\rangle_{\Sigma},
        \label{eq:stringmeasure}
\end{equation}
where
\begin{equation}
        \cB=\sum_k \dd t^k\,\cB_{t^k}.
\end{equation}
This is the form used in the external string notes.

The holomorphic reformulation is essential for the numerical approach.  If the
complex structure is described by transition functions \(z_i=f_{ij}(z_j;t)\),
then \eqref{eq:Bmetric} can be written only in terms of those transition
functions:
\begin{equation}
        \cB_{t^k}
        =
        \sum_{(ij)}
        \int_{C_{ij}}
        \left[
        \frac{\dd z_i}{2\pi\ii}\,
        b_{z_i z_i}
        \left.\frac{\partial z_i}{\partial t^k}\right|_{z_j}
        -
        \frac{\dd\bar z_i}{2\pi\ii}\,
        b_{\bar z_i\bar z_i}
        \left.\frac{\partial \bar z_i}{\partial t^k}\right|_{\bar z_j}
        \right].
        \label{eq:Btransition}
\end{equation}
This formula is the continuum reason that a construction based on sewing maps
and holomorphic data is enough to recover the string integrand.

For a genus \(g\geq2\) vacuum amplitude in complex coordinates
\(\tau^a\), \(a=1,\ldots,3g-3\), the top form is
\begin{equation}
        \Omega_{6g-6}
        =
        \left\langle
        \prod_{a=1}^{3g-3}
        \dd\tau^a\,\dd\bar\tau^a\,
        \cB_{\bar\tau^a}\cB_{\tau^a}
        \right\rangle_{\Sigma}^{\rm matter+ghost}.
        \label{eq:vacuumform}
\end{equation}
This unmarked formula is the standard continuum benchmark.  The one-face
ribbon construction used below works instead on \(\cM_{g,1}\).  For the
partition-function application we regard the marked puncture as carrying the
identity operator, so the corresponding top form has degree
\begin{equation}
        \dim_{\bR}\cM_{g,1}=6g-4.
        \label{eq:mg1dim}
\end{equation}
Equivalently, in local complex coordinates \(u^A\), \(A=1,\ldots,3g-2\), on
\(\cM_{g,1}\), one may write schematically
\begin{equation}
        \Omega_{6g-4}^{\mathbf 1}
        =
        \left\langle
        \prod_{A=1}^{3g-2}
        \dd u^A\,\dd\bar u^A\,
        \cB_{\bar u^A}\cB_{u^A}\,
        (c\widetilde c\,\mathbf 1)(p)
        \right\rangle_{\Sigma,p}^{\rm matter+ghost},
        \label{eq:identityform}
\end{equation}
where \(p\) is the marked point.  In a coordinate patch split as
\((\tau^a,p)\), the two \(B\)-insertions associated with moving \(p\) strip the
\(c\widetilde c\) factor and leave an integrated identity insertion.  This is
a marked partition-function density, not a claim that the identity insertion is
an on-shell scattering vertex.  More explicitly, the puncture-moving Beltrami
differentials reduce to the local contour modes \(b_{-1}\) and
\(\widetilde b_{-1}\) around \(p\); their action on the
\(c\widetilde c\) at the marked point supplies the local Jacobian for
integrating over the marked fiber.  This is the usual BRST coordinate-split
descent from an unintegrated marked insertion to an integrated insertion.
If period-matrix coordinates are available, the \(b\)-ghost zero modes are
paired with holomorphic quadratic differentials.  In genus two the quadratic
differentials may be written
\begin{equation}
        S_1=\omega_1^2,\qquad
        S_2=\omega_2^2,\qquad
        S_3=\omega_1\omega_2,
\end{equation}
with \(A\)-normalized holomorphic one-forms \(\omega_I\).  Writing \(V_X\) for
the noncompact target-space zero-mode volume, the genus-two critical bosonic
string vacuum form is
\begin{equation}
        \Omega_6
        =
        -
        \frac{V_X}{(4\pi^2\A)^{26}(\det\im\Omega)^{13}}
        \left|
        \frac{2\pi\ii}{\chi_{10}(\Omega)}
        \dd\Omega_{11}\dd\Omega_{22}\dd\Omega_{12}
        \right|^2,
        \label{eq:genus2critical}
\end{equation}
up to the standard overall normalization convention.  Equation
\eqref{eq:genus2critical} is the principal genus-two analytic benchmark for
the ghost and determinant machinery.

\section{One-face Strebel cells and disk sewing}

By the Strebel theorem, a point of \(\cM_{g,1}\times\bR_+\) can be represented
by a metric ribbon graph with one face.  In the generic top-dimensional cell
the graph is trivalent.  Let the oriented boundary of the unique face be traced
once.  Each edge \(e\) appears twice in this boundary word, with opposite
orientations.  The disk model places the boundary word on \(|z|=1\), with
integer edge length \(\ell_e\), so that the boundary has \(L=2\sum_e\ell_e\)
lattice sites.

Let \(S_a\), \(a=1,\ldots,2E\), be the oriented boundary segments.  If the two
occurrences of edge \(e\) are \(S_{a(e)}\) and \(S_{b(e)}\), sewing imposes an
orientation-reversing identification between the corresponding arcs of the
unit circle.  In the noncompact boson this is simply equality of the boundary
field values.  In the compact boson it is equality up to an integral shift,
which will be described in Section \ref{sec:compact}.

The boundary prevertices are the points of \(|z|=1\) where adjacent boundary
segments meet.  There are
\[
        N_{\rm pv}=2E=3V=12g-6
\]
such points in a trivalent one-face cell; triples of boundary prevertices are
identified to the \(V\) trivalent surface vertices after sewing.  If
\(q_r\), \(r=1,\ldots,N_{\rm pv}\), denote the boundary prevertices, a
holomorphic one-form on the sewn surface pulls back to the disk with cubic-root
singularities at these points.  The numerical ansatz used by the improved
higher-genus solver has the form
\begin{equation}
        \omega(z)
        =
        \left[\prod_{r=1}^{N_{\rm pv}}
        \left(1-\frac{z}{q_r}\right)^{-1/3}\right]
        P(z)\,\dd z,
        \label{eq:oneformansatz}
\end{equation}
where \(P\) is represented by a polynomial in the truncated disk basis.  The
branch choices are those inherited from evaluating the singular factor inside
the disk.
The exponent \(-1/3\) has a local geometric origin.  At a trivalent surface
vertex, the three boundary prevertices in its equivalence class contribute
three disk half-neighborhoods, with total angle \(3\pi\) before passing to a
smooth surface coordinate.  On each branch, a smooth local coordinate \(\xi\)
is therefore related to the disk coordinate by
\(\xi\sim (z-q_r)^{2/3}\), and a regular holomorphic differential \(\dd\xi\)
pulls back as \((z-q_r)^{-1/3}\dd z\) times a regular factor.

\begin{remark}
The lattice data are a regulator for the CFT on the sewn Riemann surface, not a
piecewise-linear approximation to the final worldsheet metric.  The continuum
complex structure is obtained from the sewing, while the Strebel/disk metric
only fixes a Weyl frame in which the regulator is defined.
\end{remark}

\section{Matter determinant from the identity wavefunctional}

For one free noncompact boson on the disk, the boundary field
\[
        X(\sigma)=x_0+\sum_{n\neq0}X_n\ee^{\ii n\sigma}
\]
has identity wavefunctional
\begin{equation}
        \Psi_1[X]
        \propto
        \exp\left[-\frac{1}{\A}\sum_{n\geq1}n X_nX_{-n}\right].
        \label{eq:identitywf}
\end{equation}
The scripts use a circulant lattice kernel for this quadratic form.  Up to an
overall convention-dependent factor, the kernel is
\begin{equation}
        \wh K_L(r)
        =
        \frac{\sin(\pi/L)}
        {L\left[\cos(2\pi r/L)-\cos(\pi/L)\right]},
        \qquad r\in\bZ/L\bZ.
        \label{eq:circkernel}
\end{equation}
Let \(K_L\) denote the associated circulant matrix,
\begin{equation}
        (K_L)_{ij}=\wh K_L((j-i)\bmod L),
        \label{eq:circmatrix}
\end{equation}
with indices understood modulo \(L\).  This matrix has the constant vector as
its null vector and approaches the continuum Dirichlet-to-Neumann quadratic form
in the low Fourier modes.  The files \path{partition_function.py} and
\path{compact_partition.py} differ by a factor of \(2\) in this raw matrix;
this is a regulator normalization convention, provided it is used consistently
when comparing determinants and winding forms.
With the sign convention in \eqref{eq:circkernel}, the formula agrees with
\path{direct_mat_n_fast} in \path{partition_function.py}; the corresponding
routine in \path{compact_partition.py} returns one half of the same matrix.
The positive object used in determinant computations is the sewn, reduced
quadratic form \(A'_{\Gamma,L}\), after removal of the constant zero mode.

Let \(X\in\bR^L\) be the vector of boundary values before sewing.  Sewing the
disk boundary is represented by a linear map
\[
        X=M_{\Gamma}y
\]
from independent real variables \(y\) to the full boundary vector.  The
quadratic form on the sewn surface is
\begin{equation}
        A_{\Gamma,L}=M_{\Gamma}^{T}K_LM_{\Gamma}.
        \label{eq:sewnmattermatrix}
\end{equation}
It still has one null vector, the constant mode.  If \(Q\) is any matrix whose
columns span a complement to constants, define
\begin{equation}
        A'_{\Gamma,L}=Q^{T}A_{\Gamma,L}Q.
\end{equation}
Then the regulated noncompact one-boson partition function has the form
\begin{equation}
        Z^{\rm nc}_{\Gamma,L}
        =
        C_L\,V_X\,(\det A'_{\Gamma,L})^{-1/2},
        \label{eq:noncompactdet}
\end{equation}
where \(C_L\) is a local regulator constant and \(V_X\) is the target-space
zero-mode volume.  For \(D\) noncompact bosons the determinant factor is
raised to the \(D\)-th power.

\section{Compact boson and the winding lattice}
\label{sec:compact}

Let \(X\sim X+2\pi R\).  For each edge \(e\), the two sewn boundary segments
may differ by a shift \(2\pi R s_e\) with \(s_e\in\bZ\).  Around every cubic
vertex the shifts must close.  With an orientation on the graph, this closure
condition is the integer incidence equation
\begin{equation}
        D s=0,\qquad s\in\bZ^E,
        \label{eq:incidence}
\end{equation}
where \(D\) is the \(V\times E\) oriented incidence matrix.  Since the graph is
connected,
\[
        \rank D=V-1,
        \qquad
        \dim_{\bZ}\Ker D=E-V+1=2g.
\]
Choose an integral basis
\begin{equation}
        B:\bZ^{2g}\longrightarrow \Ker_{\bZ}D,\qquad s=Bn.
        \label{eq:windingbasis}
\end{equation}
In the code this basis is constructed from graph cycles, equivalently from a
spanning tree or from the stored graph incidence data.

In a fixed shift sector, the boundary field can be written
\begin{equation}
        X=M_{\Gamma}y+2\pi R\,W s,
        \label{eq:compactaffine}
\end{equation}
where \(W\) inserts the chosen shifts along the identified sites.  Substitution
into the disk quadratic form gives
\begin{equation}
        S_L(y,s)
        =
        \frac{1}{2}y^TA'y
        +2\pi R\,y^TU s
        +(2\pi R)^2\,\frac{1}{2}s^TCs.
\end{equation}
After integrating over \(y\) on the nonzero-mode subspace, the winding sector
is weighted by the Schur complement
\begin{equation}
        T_{\rm Schur}=C-U^T(A')^{-1}U,\qquad
        T'_{\rm Schur}=B^TT_{\rm Schur}B.
        \label{eq:Tprime}
\end{equation}
The compact-boson scripts use a normalized matrix, denoted \(T'\) below and
stored as \path{T_reduced}, for which all factors from the \(2\pi R\) shifts,
the factor \(1/2\) in the Gaussian action, and the raw-kernel normalization have
already been absorbed.  Equivalently, \(T'\) is defined by the winding weight
\[
        \exp[-4\pi R^2 n^TT'n].
\]
With this convention,
\begin{equation}
        Z^{(g)}_{\Gamma,L}(R)
        =
        2\pi R\,C_L(\det A')^{-1/2}
        \Theta_{2g}(4\ii R^2T'),
        \label{eq:compactlattice}
\end{equation}
where
\begin{equation}
        \Theta_N(\Xi)
        =
        \sum_{n\in\bZ^N}
        \exp\left(\pi\ii\,n^T\Xi n\right).
\end{equation}
For \(\Xi=4\ii R^2T'\), the summand is
\(\exp[-4\pi R^2 n^TT'n]\).

The continuum compact-boson partition function on a genus \(g\) surface with
period matrix \(\Omega\), in the same \(\A=1\) convention and up to a
radius-independent factor, is
\begin{equation}
        Z_{\rm comp}^{(g)}(\Omega,R)
        =
        2\pi R\, Z_{\rm bos}^{(g)}(\Omega)
        \sum_{m,n\in\bZ^g}
        \exp\left[
        -\pi R^2
        (m+\Omega n)^T(\im\Omega)^{-1}(m+\bar\Omega n)
        \right].
        \label{eq:continuumcompact}
\end{equation}
The code checks that ratios of the lattice theta factor agree with
\eqref{eq:continuumcompact}.  For a general positive matrix, Poisson summation
does not return the same theta series; it returns the theta series of the dual
quadratic form.  Let \(N=2g\) and \(G=4T'\).  Then
\begin{equation}
        \Theta_N(\ii R^2G)
        =
        R^{-N}(\det G)^{-1/2}
        \sum_{m\in\bZ^N}
        \exp\left[-\pi R^{-2}m^TG^{-1}m\right].
        \label{eq:poissontheta}
\end{equation}
Thus a same-matrix T-duality statement is equivalent to a self-duality
condition on the normalized integral winding lattice: \(G^{-1}\) must be
integrally equivalent to \(G\), i.e. \(U^TG^{-1}U=G\) for some
\(U\in GL(N,\bZ)\).  This includes \(\det G=1\).  More generally, any residual
factor would have to be independent of both \(R\) and the moduli to be
absorbed into the overall convention.  The matrix \(T'\) stored as
\path{T_reduced} is normalized so that the following identity is the numerical
diagnostic being tested:
\begin{equation}
        R^g\Theta_{2g}(4\ii R^2T')
        =
        R^{-g}\Theta_{2g}(4\ii R^{-2}T'),
        \label{eq:tdualtheta}
\end{equation}
which implies
\begin{equation}
        Z^{(g)}(R)=R^{2-2g}Z^{(g)}(R^{-1})
\end{equation}
with the compact zero-mode factor included.  In this note \eqref{eq:tdualtheta}
is used as a numerical check of the compact-boson construction and of the
normalization of the winding lattice, not as an independent theorem about
arbitrary positive matrices \(T'\).

\section{Holomorphic data from the sewn disk}

The same ribbon data determine the period matrix and the ingredients of the
Verlinde-Verlinde \(bc\) correlator.

\subsection{Holomorphic one-forms}

The improved higher-genus solver in \path{ell_to_tau.py} uses
\eqref{eq:oneformansatz}.  Let \(m=\sum_e\ell_e\).  The regular part is
expanded as
\begin{equation}
        P_I(z)=\sum_{a=0}^{m-1}c_{I,a}z^a.
\end{equation}
For each edge, the two boundary arcs are sampled at half-integer lattice
points.  The pulled-back differential must match under orientation-reversing
sewing.  In the implemented convention the seam equation has the schematic
form
\begin{equation}
        \omega(z_R)+\omega(z_L)=0
        \label{eq:seamconstraint}
\end{equation}
at the paired sample points, after including the singular factor.  The code
constructs the matrix of these homogeneous constraints, finds its \(g\)
smallest singular directions, and then performs a Hermitian block solve to
obtain a stable basis of \(g\) one-forms.

\subsection{Homology and periods}

For a one-face graph, each edge defines a chord connecting the two occurrences
of that edge on the boundary word.  Two such chords intersect exactly when
their endpoints interleave along the boundary.  This gives an antisymmetric
integer matrix \(J\) on the free abelian group generated by edge chords.  Its
rank must be \(2g\).  The code searches for a symplectic basis
\[
        \{\alpha_I,\beta_I\}_{I=1}^g
\]
made from raw edge chords or small integer combinations, satisfying
\begin{equation}
        \alpha_I\cdot\alpha_J=0,\qquad
        \beta_I\cdot\beta_J=0,\qquad
        \alpha_I\cdot\beta_J=\delta_{IJ}.
\end{equation}

Let \(\omega_i\), \(i=1,\ldots,g\), be the unnormalized forms found above.
Using disk antiderivatives evaluated at edge midpoints, the code computes
\begin{equation}
        A_{ij}=\oint_{\alpha_j}\omega_i,
        \qquad
        B_{ij}=\oint_{\beta_j}\omega_i.
\end{equation}
The normalized period matrix is then
\begin{equation}
        \Omega=A^{-1}B.
        \label{eq:periodmatrix}
\end{equation}
The normalized forms are \(A^{-1}\omega\) in the same row convention.  Numerical
diagnostics check that \(\Omega\) is symmetric and that \(\im\Omega\) is
positive definite.

\subsection{Abel map, prime form, and sigma}

The Abel map is
\begin{equation}
        \zeta_I(z)=\int_{z_0}^{z}\wh\omega_I,
\end{equation}
where \(z_0\) is the fixed disk basepoint and \(\wh\omega_I\) are the
\(A\)-normalized holomorphic one-forms.  The Riemann theta function convention
is
\begin{equation}
        \theta[\epsilon,\delta](y|\Omega)
        =
        \sum_{n\in\bZ^g}
        \exp\left[
        \pi\ii(n+\epsilon)^T\Omega(n+\epsilon)
        +2\pi\ii(n+\epsilon)^T(y+\delta)
        \right].
\end{equation}
This is the convention implemented in the Python theta routine; other common
half-characteristic conventions differ by characteristic-dependent phase
factors.
For an odd characteristic \(\nu\), the prime form is
\begin{equation}
        E(z,w)
        =
        \frac{\theta[\nu](\zeta(z)-\zeta(w)|\Omega)}
        {\sqrt{\omega_\nu(z)\omega_\nu(w)}}.
        \label{eq:primeform}
\end{equation}
The result is independent of the odd characteristic; this is tested
numerically.

The Riemann constant vector \(\Delta\) is constructed canonically for
\(g\geq2\) using a canonical divisor from zeros of a normalized holomorphic
one-form and a half-period search selected by Riemann vanishing.  This replaces
the older direct cycle-integral formula, which is known not to be reliable for
\(g\geq2\).

The \(\sigma\) function entering the \(bc\) correlator is used through ratios
and through a normalization fixed by the \(\lambda=1\) Verlinde-Verlinde
identity.  The auxiliary divisor points used for this normalization must be
generic and disjoint from all \(b\)- and \(c\)-insertions.  Same-surface ratios
are checked to be stable under changing this auxiliary divisor.  In practice
the scripts reject choices in which divisor points collide with insertion
points or make the relevant theta, prime-form, or sigma-normalization
denominators numerically singular.

\section{The bc system and the chiral determinant}

For a \((b_\lambda,c_{1-\lambda})\) system, the selection rule is
\begin{equation}
        n_c-n_b=(1-2\lambda)(g-1).
        \label{eq:bcselection}
\end{equation}
In the conventions used by the current implementation, \(Z_{\rm ch}\) is the
primary chiral determinant factor.  The historical symbol \(Z_1\) is then
defined by
\[
        Z_1:=Z_{\rm ch}^{-2},\qquad Z_{\rm ch}=Z_1^{-1/2}.
\]
The holomorphic correlator is
\begin{equation}
\begin{aligned}
&\left\langle
\prod_{i=1}^{n_b} b_\lambda(z_i)
\prod_{j=1}^{n_c} c_{1-\lambda}(w_j)
\right\rangle_{\Sigma}
\\
&\quad =
Z_{\rm ch}\,
\theta\left(
\sum_i\zeta(z_i)-\sum_j\zeta(w_j)-(2\lambda-1)\Delta
\mid \Omega
\right)
\\
&\qquad\times
\frac{
\prod_{i<i'}E(z_i,z_{i'})
\prod_{j<j'}E(w_j,w_{j'})
\prod_i\sigma(z_i)^{2\lambda-1}}
{\prod_{i,j}E(z_i,w_j)
\prod_j\sigma(w_j)^{2\lambda-1}}.
\end{aligned}
        \label{eq:vv}
\end{equation}
The function \path{bc_correlator_geometric_factor} returns the theta,
prime-form, and sigma product in \eqref{eq:vv}; equivalently, it is the
holomorphic correlator with the single \(Z_{\rm ch}=Z_1^{-1/2}\) prefactor
stripped off.

The chiral determinant is fixed from the renormalized matter determinant.  If
\(\exp c_{\det}(\Omega)\) denotes the finite part extracted from
\(-\frac12\log\det A'\), the canonical convention in the implementation notes
is
\begin{equation}
        |Z_{\rm ch}(\Omega)|^2
        =
        [\det\im\Omega]^{1/2}\,\ee^{c_{\det}(\Omega)}
        \label{eq:zch}
\end{equation}
up to one moduli-independent normalization factor.  Equivalently,
\[
        |Z_1|^2=|Z_{\rm ch}|^{-4}.
\]
This convention is important: older text may call \(Z_1\) the chiral-boson
partition function, while the current code uses \(Z_{\rm ch}=Z_1^{-1/2}\) as
the prefactor in \eqref{eq:vv}.

\section{Large-L renormalization}

The raw lattice quantities are not finite as \(L\to\infty\).  The compact-boson
working notes introduce the counterterm scheme by fitting the self-dual compact
partition function,
\begin{equation}
        \log Z^{(g)}_{\Gamma,L}(R=1)
        =
        c_{\rm comp}(\Omega)+\gamma L+\alpha(g)\log L+o(1).
        \label{eq:compactfit}
\end{equation}
The \(bc\)/Mumford pipeline uses a determinant-only fit:
\begin{equation}
        -\frac{1}{2}\log\det A'_{\Gamma,L}
        =
        c_{\det}(\Omega)+\gamma_{\det} L+\alpha_{\det}\log L+o(1),
        \label{eq:detfit}
\end{equation}
where the code fits \(\gamma_{\det}\) and \(\alpha_{\det}\) directly for the
genus and graph families used in the ratio check.  Thus \(\ee^{c_{\det}}\), not
the compact theta sum, is the finite determinant factor used below in
\eqref{eq:zch}.

The empirical self-dual compact-partition values quoted in \path{Strebel.tex}
and in the string Monte Carlo note are
\[
        \gamma\simeq \frac{1}{4}\log2,\qquad
        \alpha_0\simeq -0.86112,\qquad
        \alpha_1\simeq -0.27778.
\]
Here \(\alpha(g)=\alpha_0+\alpha_1g\) refers to the compact self-dual fit
\eqref{eq:compactfit}.  The actual Mumford checks do not assume the quoted
\(\alpha_0,\alpha_1\): they fit the determinant coefficient
\(\alpha_{\det}\) from \(A'\)-data in the selected genus.
No analytic derivation of these constants was found in the local notes or
scripts.  Until a heat-kernel, Liouville, or regulator calculation is supplied,
the decimal fit coefficients should be treated as empirical inputs to this
counterterm scheme.
The script \path{genus1_to_4_scaling_sweep.py} is the current reproduction
driver for the compact large-\(L\) fits; for the self-dual coefficients above one
should run the sweep at \(R=1\) and record the chosen \(L\)-window, graph
families, and theta cutoffs.
With the current script defaults for the \(L\)-window, an explicit command is
\begin{quote}
\small
\begin{verbatim}
./.venv/bin/python "covariant formalism/python/genus1_to_4_scaling_sweep.py" \
  --R 1.0 --L-start 2000 --L-stop 5000 --L-step 200 \
  --json-out /tmp/radius_1_scaling.json
\end{verbatim}
\end{quote}
The exact status of these constants is an open derivation problem for a final
draft; they are not used below as an independent theorem.
The corresponding compact-fit local counterterm convention is
\begin{equation}
        S_{\rm ct}
        =
        \gamma L+(\alpha_0+\alpha_1)\log L
        -
        \frac{\alpha_1}{8\pi}\log L
        \int_{\Sigma}\dd^2\sigma\,\sqrt g\,R(g).
        \label{eq:counterterm}
\end{equation}
Since
\[
        \int_{\Sigma}\sqrt g\,R(g)=4\pi(2-2g),
\]
\eqref{eq:counterterm} reproduces \(\alpha_0+\alpha_1g\) in
\eqref{eq:compactfit}.  In the determinant-only scheme used for
\eqref{eq:zch}, the fitted finite determinant factor is
\begin{equation}
        [(\det A')^{-1/2}]_{\rm ren}=\ee^{c_{\det}(\Omega)}.
\end{equation}

This is a regulator definition, not a proof of universal convergence.  The
evidence for the compact fit is that the same \(\gamma\) and genus-linear
\(\alpha(g)\) stabilize across graph families, and that the resulting finite
part satisfies the genus-one Weyl-anomaly comparison with the analytic
compact-boson partition function.  The evidence for the determinant fit is the
Mumford-form agreement after \(\gamma_{\det},\alpha_{\det}\) are fit from
same-genus \(A'\)-families.

\section{Numerical bosonic string integrand}

The CFT part of the critical bosonic string measure can be represented in
terms of a Mumford form.  Let \(q_a(z)\), \(a=1,\ldots,3g-3\), be a basis of
holomorphic quadratic differentials.  For generic points \(z_b\), define
\begin{equation}
        D_q(z_1,\ldots,z_{3g-3})=\det[q_a(z_b)].
\end{equation}
The numerical \(bc\) correlator and the renormalized determinant determine the
unmarked local Mumford factor through
\begin{equation}
        \left|
        \left\langle \prod_{b=1}^{3g-3} b(z_b)\right\rangle_{\Sigma}
        \right|^2
        |Z_{\rm ch}(\Omega)|^{52}
        =
        C_g\,|D_q(z_1,\ldots,z_{3g-3})|^2\,|\mu_g(\Omega)|^2.
        \label{eq:mumfordpoint}
\end{equation}
Here \(C_g\) is moduli independent, and the bracket denotes the full ghost
correlator in the convention of \eqref{eq:vv}, including its single
\(Z_{\rm ch}\) prefactor.  The additional factor \(|Z_{\rm ch}|^{52}\) in
\eqref{eq:mumfordpoint} is the matter contribution of the \(26\) critical
bosons.  If one instead used the stripped geometric factor returned by
\path{bc_correlator_geometric_factor}, the displayed power of
\(|Z_{\rm ch}|\) would have to be shifted accordingly.  Equation
\eqref{eq:mumfordpoint} is the point-insertion version of the same section that
appears after contracting \(b\)-ghost zero modes with Beltrami differentials in
the string measure.
In the coordinate-split description of the identity-marked \(\cM_{g,1}\)
density, the unmarked Mumford factor above is the CFT scalar factor pulled back
along the forgetful map to the underlying surface.  The remaining marked
coordinate is accounted for by the \(B\)-insertions and \(c\widetilde c\)
factor in \eqref{eq:identityform}.

For genus two,
\begin{equation}
        \mu_2(\Omega)=\frac{2\pi\ii}{\chi_{10}(\Omega)}.
        \label{eq:mu2}
\end{equation}
The implemented ratio check is precisely
\begin{equation}
        \frac{|\vev{bbb}_1|^2}{|\vev{bbb}_2|^2}
        \left(
        \frac{|Z_{\rm ch,1}|^2}{|Z_{\rm ch,2}|^2}
        \right)^{26}
        =
        \frac{|D_{q,1}|^2}{|D_{q,2}|^2}
        \frac{|\chi_{10}(\Omega_2)|^2}{|\chi_{10}(\Omega_1)|^2},
        \label{eq:genus2check}
\end{equation}
up to numerical error.

For genus three, the six quadratic differentials are
\[
        \omega_1^2,\quad \omega_1\omega_2,\quad \omega_1\omega_3,\quad
        \omega_2^2,\quad \omega_2\omega_3,\quad \omega_3^2.
\]
For nonhyperelliptic genus-three surfaces,
\begin{equation}
        |\mu_3(\Omega)|^2\propto \frac{1}{|\chi_{18}(\Omega)|},
        \qquad
        \chi_{18}=\prod_{\delta\ {\rm even}}\theta[\delta](0|\Omega),
        \label{eq:mu3}
\end{equation}
where the product is over the \(36\) even spin structures.  The current
genus-three check compares the two sides of \eqref{eq:mumfordpoint} using
\eqref{eq:mu3}; the documented relative errors are at the \(10^{-4}\) to
\(10^{-3}\) level for the tested examples.

For genus four and above, the products \(\omega_I\omega_J\) are not an
independent basis without relations.  On the generic nonhyperelliptic genus-four
locus there is one quadric relation among the ten products.  The script
\path{genus4_bc_local_mumford_check.py} therefore performs a same-surface
local-chart check.  On a fixed surface, choose two generic nine-tuples of
\(b\)-insertion points \(P_1\) and \(P_2\).  For each usable chart obtained by
omitting a quadratic monomial with nonzero quadric coefficient, form the two
determinants \(D_q^{(r)}(P_1)\) and \(D_q^{(r)}(P_2)\).  The diagnostic compares
\[
        \frac{|\vev{b^9(P_1)}|^2}{|\vev{b^9(P_2)}|^2}
        \quad\hbox{with}\quad
        \frac{|D_q^{(r)}(P_1)|^2}{|D_q^{(r)}(P_2)|^2}
\]
on that same surface, and reports the relative error, the quadric residual, and
the spread across usable omission charts.  Hyperelliptic or otherwise special
loci require separate treatment.

\subsection{Compact radius dependence}

For a \(c=1\) matter factor consisting of a compact boson of radius \(R\) and a
radius-independent compensating CFT, the radius dependence of the fixed-surface
integrand is
\begin{equation}
        \cR_R(\Sigma)
        =
        \frac{Z_{\rm comp}^{\rm ren}(\Sigma,R)}
        {Z_{\rm comp}^{\rm ren}(\Sigma,R_0)}.
        \label{eq:radiusratio}
\end{equation}
At fixed surface, Weyl-frame and determinant factors cancel in this ratio
except for the compact winding theta functions and the zero-mode factor.  This
is the cleanest numerical observable for testing the compact-boson part of the
integrand.  The self-dual radius \(R_0=1\) is the most common reference in the
local notes.

\subsection{Cell-coordinate measure}

Let \(\ell_a\), \(a=1,\ldots,E-1=6g-4\), be local real edge-length
coordinates in a fixed-perimeter one-face ribbon cell.  The pullback of the
marked identity form \eqref{eq:identityform} to the cell is schematically
\begin{equation}
        \Phi_{\Gamma}^{*}\Omega_{6g-4}^{\mathbf 1}
        =
        \cI_{\Gamma}^{\mathbf 1}(\ell)\,
        \dd\ell_1\cdots \dd\ell_{6g-4},
        \label{eq:cellmeasure}
\end{equation}
after imposing the fixed-perimeter condition.  The scalar
\(\cI_{\Gamma}^{\mathbf 1}\) contains the matter determinant, compact winding
sum if present, and the ghost/Mumford factor appropriate to the identity-marked
surface.  A global Monte Carlo integral over \(\cM_{g,1}\) must also include
the cell automorphism factor \(1/|\Aut\Gamma|\) and a consistent treatment of
cell boundaries.
The present ratio checks do not test the automorphism factor: they compare
fixed-surface or same-cell scalar quantities, so \(1/|\Aut\Gamma|\) cancels.
The automorphism weight becomes a genuine input only in a global sum or
integration over ribbon cells.

The current code has strong tests for the CFT scalar factors and local
Mumford-section relations.  A final production integrand should still make
explicit the fixed-perimeter cell measure and the normalization convention for
the identity-marked puncture.  A later comparison with an unmarked integral on
\(\cM_g\) would require an additional pushforward or fiber-normalization step,
but that is not part of the native ribbon-cell prescription used here.

\section{Algorithmic summary}

For a fixed one-face ribbon graph \(\Gamma\) and integer edge lengths
\(\ell_e\), the implemented pipeline is:

\begin{enumerate}
\item Trace the unique boundary word and build the oriented segment pairings.
\item Construct the matter sewing matrix \(A'_{\Gamma,L}\) from the disk
identity wavefunctional.
\item For a compact boson, build the incidence matrix \(D\), an integral basis
\(B\) of \(\Ker D\), the Schur complement and script-normalized matrix \(T'\),
and the theta function \(\Theta_{2g}(4\ii R^2T')\).
\item Reconstruct \(g\) holomorphic one-forms using the cubic-prevertex
singular ansatz and the sewn-boundary matching equations.
\item Construct a symplectic homology basis from edge-chord intersections and
compute the period matrix \(\Omega=A^{-1}B\).
\item Build Abel maps, prime forms, canonical Riemann constants, and
\(\sigma\)-values.
\item Evaluate the Verlinde-Verlinde \(bc\) correlator and insert the
renormalized chiral factor \(Z_{\rm ch}\).
\item Combine the ghost correlator with the renormalized matter determinant to
obtain the local Mumford/integrand factor, comparing to analytic forms when
available.
\end{enumerate}

\begin{center}
\begin{tabular}{L{0.36\textwidth}L{0.56\textwidth}}
\hline
File & Role \\
\hline
\path{Fast_Ribbon_Generator.py} &
Enumeration of one-face ribbon-graph topologies, graph-isomorphism reduction,
automorphism data, and genus-four topology generation used by the current
diagnostics. \\
\path{partition_function.py} &
Generic one-face sewing matrices, reduced matter determinant, and ghost
wavefunctional experiments. \\
\path{compact_partition.py} &
Compact-boson winding matrices, stored genus-one/genus-two graph data, and
theta-function tests. \\
\path{ell_to_tau.py} &
Improved holomorphic one-form solver, homology basis, period matrix, local
prevertex factors, and genus-two modular forms. \\
\path{riemann_surface_tools.py} &
Surface data object, Abel maps, prime form, canonical Riemann constants,
\(\sigma\), \(bc\) correlators, determinant fits, and chiral normalization. \\
\path{genus2_bc_igusa_ratio_check.py} &
Genus-two check against the Igusa cusp form. \\
\path{genus3_bc_chi18_ratio_check.py} &
Genus-three check against the \(\chi_{18}\) relation. \\
\path{genus4_bc_local_mumford_check.py} &
Genus-four same-surface local Mumford chart check using the quadric relation. \\
\hline
\end{tabular}
\end{center}

\section{Checks and evidence}

\noindent\textbf{Check.}
The strongest existing checks are the following.  Their accuracy is limited by
different effects in different rows: finite \(L\) extrapolation and fit-window
choice for determinant renormalization, theta-series truncation, numerical
quadrature near singular prevertices, and conditioning of local determinant
charts.

\begin{center}
\begin{tabular}{L{0.31\textwidth}L{0.61\textwidth}}
\hline
Source & Reproducibility data quoted in the notes \\
\hline
\path{Strebel.tex}, \path{string MC.tex}, and
\path{genus1_to_4_scaling_sweep.py} &
At \(R=1\), the working compact partition-function fit targets
\(\gamma\simeq \frac14\log2\) and
\(\alpha(g)\simeq -0.86112-0.27778g\).  No analytic derivation of these
constants is contained in the local notes or scripts.  The sweep script records
the \(L\)-window, selected graph families, and theta cutoffs; run it with
\path{--R 1.0} for the self-dual scan.
\\
\path{covariant formalism/markdown/riemann_surface_tools.md} &
For Abel-jump consistency on stored genus-two topology \(1\), the reported
worst componentwise errors are \(6.14\times10^{-6}\), \(1.63\times10^{-6}\),
and \(7.42\times10^{-7}\) for edge lists \([50]^9\), \([100]^9\), and
\([150]^9\), respectively. \\
\path{genus2_bc_igusa_ratio_check.py} &
For \(\ell_1=[300]^9\) and
\(\ell_2=(250,250,250,250,250,250,250,250,700)\), the documented Igusa
comparison gives \( {\rm LHS}=7.139760447638384\),
\({\rm RHS}=7.139757771156295\), relative difference
\(3.75\times10^{-7}\). \\
\path{genus3_bc_chi18_ratio_check.py} and
\path{covariant formalism/markdown/genus3_bc_chi18_ratio_check.md} &
With topology \(1\), safe divisor points, and \(n_{\max}=4\), the documented
\(\chi_{18}\) comparisons against the symmetric modulus \(M_1=[300]^{15}\)
give relative errors \(4.91\times10^{-4}\), \(5.96\times10^{-4}\), and
\(4.24\times10^{-4}\) for the three listed comparison moduli. \\
\path{genus4_bc_local_mumford_check.py} &
Reports a same-surface nine-\(b\) ratio, the local quadric-chart determinant
ratio, the relative error, the quadric residual, and chart/divisor stability
diagnostics.  No representative numerical tolerance is fixed in this note; the
claim supported here is local and generic, not a global genus-four integral. \\
\hline
\end{tabular}
\end{center}

\begin{description}
\item[Genus one compact boson.]
The lattice compact partition function agrees with the analytic torus lattice
sum after mapping edge lengths to \(\tau\), and satisfies T-duality in the form
\eqref{eq:tdualtheta}.

\item[Higher-genus compact boson.]
The incidence-kernel construction gives a rank-\(2g\) winding lattice.  The
genus-\(g\) T-duality identity is then a numerical diagnostic: the notes report
tests for all stored genus-two cases, the stored genus-three cases, and one
genus-four topology.

\item[Period matrices.]
The reconstructed \(\Omega\) is symmetric up to numerical error and has
positive definite imaginary part.  Abel jumps across sewn boundary pairs land
on the period lattice with errors decreasing as the edge lengths increase.

\item[Prime form.]
Spin-structure independence of the prime form is checked numerically.  The
genus-two notes quote relative differences of order \(10^{-6}\) or smaller in
representative large-edge tests.

\item[Genus two Mumford form.]
The Igusa-ratio check \eqref{eq:genus2check} agrees at the level of roughly
\(4\times10^{-7}\) in the documented run.

\item[Genus three Mumford form.]
The \(\chi_{18}\) check uses the corrected squared-absolute relation
\[
        |\mu_3|^2\propto \frac{1}{|\chi_{18}|}.
\]
The documented runs agree at roughly \(10^{-4}\) to \(10^{-3}\).  The weaker
precision compared with the genus-two Igusa check is consistent with the
higher-dimensional theta sums, the sigma-normalization sensitivity, and the
larger number of \(b\)-ghost insertions; it should be treated as a numerical
accuracy budget to be reported, not as a new analytic assumption.

\item[Genus four local chart.]
The same-surface nine-\(b\) ratio is implemented as a comparison against the
determinant ratio in a local quadric chart, with chart spread used as a
diagnostic of numerical stability.  A draft-ready statement should quote a
specific run, relative error, quadric residual, and divisor/chart stability
tolerance.
\end{description}

\section{Reference anchors}

This note is primarily a code-to-continuum synthesis.  A polished draft should
replace the path-style references above by ordinary citations.  The minimum
source list should include the following anchors.

\begin{itemize}
\item Strebel's theory of quadratic differentials for the correspondence
between metric ribbon graphs and punctured Riemann surfaces with perimeters.
\item Kontsevich's matrix-Airy construction, and the Mulase--Penkava
exposition, for the ribbon-cell decomposition and its relation to moduli-space
forms.
\item Faddeev--Popov, Friedan--Martinec--Shenker, Belavin--Knizhnik, and
Polchinski for the covariant string measure, \(bc\)-ghost insertions, and the
BRST descent between unintegrated and integrated vertex insertions.
\item Mumford's determinant-line formalism, together with the Verlinde--
Verlinde bosonization formulae, for the chiral measure and the ghost
correlator used in the Mumford-form checks.
\item Igusa's work on Siegel modular forms for the genus-two \(\chi_{10}\)
and genus-three \(\chi_{18}\) modular forms used as analytic benchmarks.
\item Deconinck--Patterson--Swierczewski for the canonical-divisor algorithm
used in the current Riemann-constant computation.
\item Bacher--Vdovina and the local ribbon-graph enumeration scripts for
one-vertex triangulation counts and topology bookkeeping.
\end{itemize}

\section{Open points for the draft}

The following issues should be stated explicitly in any polished draft.

\begin{enumerate}
\item The one-face Strebel coordinates naturally describe
\(\cM_{g,1}\times\bR_+\).  This note works natively on fixed-perimeter
\(\cM_{g,1}\), with the marked puncture carrying the identity operator for
partition-function applications.  A comparison to an unmarked amplitude on
\(\cM_g\) would require an additional pushforward or fiber-normalization
convention.
\item The absolute normalization of the full amplitude is not fixed by the
current determinant and ghost checks.  Ratios and local Mumford sections are
the controlled quantities.
\item The determinant finite part \(\ee^{c_{\det}(\Omega)}\) is a renormalized
quantity defined by the large-\(L\) fit.  The fit should be treated as part of
the regulator definition, with convergence tests reported for every production
regime.
\item Numerical values quoted in the draft should be tied to a script, command
line, graph or topology choice, edge-length data, theta cutoff, fit window, and
code revision.  The small source table in Section~12 is a starting point, not a
replacement for a reproducibility table in a paper.
\item Lower-dimensional ribbon cells, automorphism factors, the fixed-perimeter
cell measure, and boundary behavior of the integrand are required for a global
Monte Carlo integral on \(\cM_{g,1}\).
\item The legacy one-form and Riemann-constant routines should not be used in
the final pipeline.  The improved cubic-singularity solver and canonical
Riemann-constant construction are the versions supported by the current tests.
\item In genus four and higher, the quadratic-differential basis must account
for canonical-curve relations on the relevant stratum.  The genus-two and
genus-three closed modular forms do not generalize naively.
\item The path-style references in this note should be converted into precise
\verb|\cite| entries in a polished draft, using the anchors in the preceding
section as a starting point.  When older notes disagree with current scripts or
Markdown diagnostics, the current code and the checked numerical diagnostic
should be treated as the source of record until the disagreement is resolved
explicitly.
\end{enumerate}

\end{document}
